% ============================================================
%  preamble.tex — shared preamble for the algorithm reference.
%  Typographic conventions inspired by the canonical algorithms
%  textbook: Computer Modern, two-column-friendly math, boxed
%  pseudocode with line numbers, margin-callout invariants.
% ============================================================

\usepackage[T1]{fontenc}
\usepackage[utf8]{inputenc}
\usepackage{lmodern}
\usepackage{textcomp}
\usepackage[margin=1.25in]{geometry}
\usepackage{amsmath,amssymb,amsthm}
\usepackage{mathtools}
\usepackage{graphicx}
\usepackage{xcolor}
\usepackage{microtype}
\usepackage{enumitem}
\usepackage{booktabs}
\usepackage{tabularx}
\usepackage{array}
\usepackage{caption}
\usepackage[nottoc,notlot,notlof]{tocbibind}
\usepackage{titlesec}
\usepackage{fancyhdr}
\usepackage{framed}
\usepackage[ruled,linesnumbered,noend]{algorithm2e}
\usepackage{hyperref}
\hypersetup{
    colorlinks=true,
    linkcolor=black,
    citecolor=black,
    urlcolor=black,
    pdftitle={GraphHunter Algorithms Reference},
    pdfauthor={GraphHunter Team}
}

% ------------------------------------------------------------
% Page style — running header with chapter + section, as in
% the reference algorithms text.
% ------------------------------------------------------------
\pagestyle{fancy}
\fancyhf{}
\renewcommand{\headrulewidth}{0.4pt}
\fancyhead[LE,RO]{\thepage}
\fancyhead[LO]{\small\itshape\nouppercase{\rightmark}}
\fancyhead[RE]{\small\itshape\nouppercase{\leftmark}}

% ------------------------------------------------------------
% Chapter / section styling — restrained, serif, no color.
% ------------------------------------------------------------
\titleformat{\chapter}[display]
  {\normalfont\Huge\bfseries}
  {\chaptertitlename\ \thechapter}{20pt}{\Huge}
\titleformat*{\section}{\Large\bfseries}
\titleformat*{\subsection}{\large\bfseries}
\titleformat*{\subsubsection}{\normalsize\bfseries\itshape}

% ------------------------------------------------------------
% Theorem-like environments with custom styles matching the
% "Theorem / Lemma / Invariant" layout.
% ------------------------------------------------------------
\newtheoremstyle{ghtheorem}
  {12pt}{12pt}{\itshape}{}{\bfseries}{.}{.5em}
  {\thmname{#1}\thmnumber{ #2}\thmnote{ (#3)}}
\newtheoremstyle{ghdefinition}
  {12pt}{12pt}{}{}{\bfseries}{.}{.5em}
  {\thmname{#1}\thmnumber{ #2}\thmnote{ (#3)}}

\theoremstyle{ghtheorem}
\newtheorem{theorem}{Theorem}[chapter]
\newtheorem{lemma}[theorem]{Lemma}
\newtheorem{corollary}[theorem]{Corollary}
\newtheorem{invariant}[theorem]{Invariant}

\theoremstyle{ghdefinition}
\newtheorem{definition}[theorem]{Definition}
\newtheorem{problem}[theorem]{Problem}
\newtheorem{remark}[theorem]{Remark}

% ------------------------------------------------------------
% Algorithm2e tuning: numbered lines, bold keywords, en-dash
% comments — matches the pseudocode convention of the canonical
% text. Keywords small-caps, variable names italic math.
% ------------------------------------------------------------
\SetKwSty{textsc}
\SetFuncSty{textsc}
\SetArgSty{textit}
\SetKwInOut{Input}{Input}
\SetKwInOut{Output}{Output}
\SetKwInOut{Invariant}{Invariant}
\SetKwFor{For}{for}{do}{end for}
\SetKwFor{While}{while}{do}{end while}
\SetKw{KwReturn}{return}
\SetKw{KwAnd}{and}
\SetKw{KwOr}{or}
\SetKw{KwNot}{not}
\SetKwComment{Comment}{\hfill$\triangleright$\ }{}
\DontPrintSemicolon
\LinesNumbered

% ------------------------------------------------------------
% Custom math operators and shorthands used across chapters.
% ------------------------------------------------------------
\DeclareMathOperator*{\argmax}{arg\,max}
\DeclareMathOperator*{\argmin}{arg\,min}
\DeclareMathOperator{\freq}{freq}
\DeclareMathOperator{\degout}{deg}
\newcommand{\bigO}{\mathcal{O}}
\newcommand{\N}{\mathbb{N}}
\newcommand{\Z}{\mathbb{Z}}
\newcommand{\R}{\mathbb{R}}
\newcommand{\eps}{\varepsilon}
\newcommand{\set}[1]{\{#1\}}
\newcommand{\abs}[1]{\lvert #1 \rvert}
\newcommand{\card}[1]{\lvert #1 \rvert}
\newcommand{\code}[1]{\texttt{#1}}
\newcommand{\unit}[1]{\ensuremath{\,\mathrm{#1}}}

% ------------------------------------------------------------
% "Implementation note" side-block: a lightly framed aside that
% connects the formal specification to the code that realizes it.
% ------------------------------------------------------------
\newenvironment{implnote}[1][Implementation note]{%
  \par\medskip\noindent
  \begin{framed}\small
  \textbf{#1.}\ \ignorespaces
}{\end{framed}\par\medskip}

% End of preamble.
